\section{Gaussian Mean--Variance Derivations}
\label{app:gaussian}

% MANUSCRIPT-NODE: APP-GAUSS-P01
For a scalar Gaussian with $\xi=\log\sigma$ and $z=(a-\mu)/\sigma$,
\begin{equation}
\partial_\mu\log\pi=\frac{a-\mu}{\sigma^2},\qquad
\partial_\xi\log\pi=z^2-1.
\end{equation}
A negative advantage reverses both score directions. It always pushes the mean away from the action, but its effect on scale depends on $|z|$: a near negative action with $|z|<1$ increases $\sigma$, whereas a far negative action with $|z|>1$ decreases $\sigma$. Positive updates have the opposite signs. A full Gaussian terminal state must satisfy both mean and variance equations; a stationary mean with continuing covariance movement is not a complete policy equilibrium.
% END-MANUSCRIPT-NODE: APP-GAUSS-P01

\subsection{Population mean--variance fixed point}
With $\xi=\log\sigma$, positive mass $p$, effective negative mass $q$, conditional moments $(m_+,v_+)$ and $(m_-,v_-)$,
\begin{align}
\dot\mu &= \frac{p(m_+-\mu)-q(m_--\mu)}{\sigma^2},\\
\dot\xi &= \frac{p[v_++(\mu-m_+)^2]-q[v_-+(\mu-m_-)^2]}{\sigma^2}-(p-q).
\end{align}
For $p>q$, the mean candidate is $\mu^\star=(pm_+-qm_-)/(p-q)$. Substitution into the scale equation gives
\begin{equation}
(\sigma^2)^\star=\frac{pM_+(\mu^\star)-qM_-(\mu^\star)}{p-q},
\end{equation}
where $M_\pm(\mu)=v_\pm+(\mu-m_\pm)^2$. A finite joint equilibrium therefore requires both $p>q$ and a positive numerator. This second condition is stricter in the far field because $M_-$ contains squared displacement.

